\section{Inference to the best explanation: formal separating cases}
\label{sec:annex-p16-inference-to-best-explanation}

These are full proofs of the expressly added finite illustrations in
\cref{sec:p16-inference-to-best-explanation}. They establish neither a
general definition of explanatory quality nor Lipton's philosophical
reliability claims. Those arguments remain in the attributed body.

\begin{annexitem}{annex:p16.ibe-printer}{thm:p16-ibe-printer}{All specified assignments and both replacements}
\begin{proof}
In context \(u\), the awake equation gives \(A=1\). The output
equation has three positive inputs and \(\neg J=0\), so \(O=0\).
In context \(v\), awake state is again one and all four conjuncts of
the output equation are one, so \(O=1\).

For replacement of \(J\) in \(u\), set its assigned value to zero
and keep \(P,R,T\) at one. The same total output function gives one;
the awake function, which uses only \(P\), also remains one. For
replacement of \(P\), set it to zero, retaining all other values.
Then awake state is zero and the output conjunction is zero regardless
of the other conjuncts. The original power value one is shared by
\(u,v\), so it supplies no differing coordinate in their specified
contrast. This is an exhaustive evaluation of the two contexts and two
replacements, not a causal conclusion from an unspecified observed table.
\end{proof}
\end{annexitem}

\begin{annexitem}{annex:p16.ibe-raven-logic}{thm:p16-ibe-raven-logic}{Equivalence does not imply relevance or universality}
\begin{proof}
At any object, \(R\to B\) is false exactly when \(R\) is true
and \(B\) false. The conditional \(\neg B\to\neg R\) is false
under exactly those same truth values. They agree at every object;
universal quantification over the same domain therefore preserves their
equivalence in both directions.

At a \(\neg R\land\neg B\) object the first conditional has false
antecedent and the second true consequent, so both are true. For a domain
with at least two objects choose different objects \(a,b\). Make
\(R(a),B(a)\) both false and make \(R(b)\) true, \(B(b)\) false.
Give any remaining objects arbitrary predicate values. The observed fact
at \(a\) is true and both universals false at \(b\). Thus that fact
does not entail either universal. On a one-object domain with that
observation the universals are true vacuously; this is why the countermodel
claim includes the two-object boundary. No inductive support relation has
been introduced by the logical equivalence proof.
\end{proof}
\end{annexitem}

\begin{annexitem}{annex:p16.ibe-raven-sampling}{thm:p16-ibe-raven-sampling}{Both likelihoods and their positive normalizers}
\begin{proof}
Under U, precisely 98 of 100 possible selected objects are white
nonravens under each hypothesis. Hence the evidence has likelihood
\(98/100\) under both. Its unconditional probability is
\((1/2)(98/100)+(1/2)(98/100)=98/100>0\). The joint probability
of \(H_0\) and that evidence is \((1/2)(98/100)\); division
gives posterior \(1/2\).

Under V, a black raven has likelihood one under \(H_0\) and
\(1/2\) under \(H_1\). Its unconditional probability is
\((1/2)1+(1/2)(1/2)=3/4>0\). Its joint probability with
\(H_0\) is \(1/2\), so the posterior is \((1/2)/(3/4)=2/3\).
The posterior of \(H_1\) is the remaining \(1/3\), positive;
thus the observation does not eliminate the counterhypothesis. Both
calculations use the declared protocol and equal priors; neither follows
from the grammar of a positive or contrapositive instance alone.
\end{proof}
\end{annexitem}

\begin{annexitem}{annex:p16.ibe-bayes-scope}{thm:p16-ibe-bayes-scope}{A common normalization and the considered-set conditioning}
\begin{proof}
Positive \(P(B_0)\) makes conditioning on background legitimate.
Since the \(H_i\) partition the space, finite total probability gives
\[
 P(E\mid B_0)=\sum_iP(H_i\mid B_0)P(E\mid H_i,B_0)=Z>0.
\]
The positive prior \(\pi_i\) also makes every \((H_i,B_0)\)
conditioning event positive. Conditional multiplication gives
\(P(H_i,E\mid B_0)=\pi_i\ell_i\); dividing by \(Z\)
proves \(a_i=\pi_i\ell_i/Z\). If \(\ell_j>0\), then
\(a_j>0\), so dividing the two formulas cancels \(Z\) and
proves the odds identity. A zero numerator likelihood gives odds zero;
no ratio is asserted when its denominator likelihood is zero.

For the considered event, disjointness gives
\(P(C_I\mid E,B_0)=\sum_{i\in I}a_i=r>0\). For \(i\in I\),
\(H_i\cap C_I=H_i\), so a further legal conditioning gives
\(q_i=a_i/r\), equivalently \(a_i=rq_i\). The stated triple
has sum one and nonnegative entries. For \(I=\{1,2\}\),
\(r=.20\) and \(q_1=.18/.20=.90\).

More generally fix any interior \(q_1\) and any \(0<r<1\), with
at least one other considered index and one outside index. Assign posterior
\(rq_1\) to the first, \(r(1-q_1)\) to the second, and \(1-r\)
to the outside event. These are positive masses summing to one and can
be realized with evidence certain. Letting \(r\) decrease makes the
absolute first mass arbitrarily small while retaining its considered-set
rank whenever \(q_1>1/2\). This proves the no-lower-bound statement
for these incomplete considered sets. For a singleton considered set,
assign it mass \(r\) and an outside event mass \(1-r\). Its
considered-set posterior is one while its absolute posterior is \(r\),
again arbitrarily small. More generally, multiplying any positive
considered-set probability vector by \(r\) preserves every conditional
rank while assigning the remaining \(1-r\) to an outside event.
If the set is exhaustive, \(r=1\)
and this particular problem disappears; exhaustion itself needs justification.
\end{proof}
\end{annexitem}

\begin{annexitem}{annex:p16.ibe-bonus}{thm:p16-ibe-bonus}{Exactly when the bonus leaves the posterior unchanged}
\begin{proof}
Put \(D=\sum_jb_j a_j\). Positive finite weights and positive
posterior masses give \(D>0\). If \(a_i^*=a_i\) for all \(i\),
then \(b_i a_i/D=a_i\). Division by positive \(a_i\) gives
\(b_i=D\) at every index, so the weights are constant. Conversely,
if every \(b_i=c>0\), normalization of the posterior gives
\(D=c\sum_i a_i=c\), and the alternative mass is \(a_i\).
The proof concerns the positive posterior support. Weights at zero-mass
events would not be constrained by this equality, which is why the statement
uses positive masses throughout. No philosophical premise equating
explanatory virtue with an independently warranted probability has been proved.
\end{proof}
\end{annexitem}

\begin{annexitem}{annex:p16.ibe-contradictories}{thm:p16-ibe-contradictories}{The exact probabilistic content of two ranking premises}
\begin{proof}
Complementation gives \(P(\neg H)=1-P(H)\). Thus
\(P(H)>P(\neg H)\) is exactly \(P(H)>1-P(H)\), or
\(2P(H)>1\), proving the first implication. If the two probabilities
are equal, the same identity gives \(P(H)=1/2\) and no strict conclusion.

For \(K\subseteq\neg H\), disjoint additivity yields
\(P(\neg H)=P(K)+P(\neg H\setminus K)\ge P(K)\).
With the stipulated \(P(K)>P(H)\), transitivity gives
\(P(\neg H)>P(H)\). The reverse relation follows only if the
appropriate comparison premises are supplied. An unexplained rank, or
an assertion that some hypothesis is lovely, is not any of these numerical
premises and therefore is not covered by the proof.
\end{proof}
\end{annexitem}

\begin{annexitem}{annex:p16.ibe-selection-fit}{thm:p16-ibe-selection-fit}{Fixed and data-dependent matching are different events}
\begin{proof}
There are \(2^k\) binary strings of length \(k\), all equally likely.
Exactly one equals the beforehand fixed string \(s\), so
\(P(X=s)=1/2^k\). In the second procedure \(S(X)=X\) for every
possible observation. Consequently the event \(S(X)=X\) is the whole
outcome space and has probability one. At \(k=8\), the count is 256.

The selected-string equality is not the event that a fixed \(s\) matched
without selection. The proof depends on uniformity and on searching all
strings with the stipulated selection map; it does not apply to every
scientific theory-construction process. In particular it does not compute
either procedure's probability of selecting a true causal account.
\end{proof}
\end{annexitem}
